% !PRINTED 525
\weberpage{525}
%<<CONT
that everything can easily be surveyed, and yet the most important phenomena of group formation come to light in the process.

From four figures, 0, 1, 2, 3, one can form 24 permutations, and so high, therefore, is the degree of the \gesperrt{Galois} group of the \gesperrt{general} equation of the fourth degree in the field of the rational functions of the coefficients (cf. the end of §. 149). We represent these permutations by their cycles in the following way, where 1 denotes the identical permutation:
\weq{1}{\begin{gathered}
P = 1, \quad (0,1), \quad (0,2), \quad (0,3), \quad (1,2), \quad (1,3), \quad (2,3) \\
(0,1)(2,3), \quad (0,2)(1,3), \quad (0,3)(1,2) \\
(0,1,2), \quad (0,1,3), \quad (0,2,3), \quad (1,2,3) \\
(0,2,1), \quad (0,3,1), \quad (0,3,2), \quad (1,3,2) \\
(0,1,2,3), \quad (0,1,3,2), \quad (0,2,3,1) \\
(0,2,1,3), \quad (0,3,1,2), \quad (0,3,2,1).
\end{gathered}}

We thus have in $P$, besides the identical permutation, six transpositions, eight three-term cycles, six four-term cycles, and three permutations of the form $(0,1)(2,3)$, which we shall call \gesperrt{transposition pairs}.

In $P$ the alternating group of degree 12 is contained as a normal divisor. This group can contain no single transposition and no four-term cyclic permutation, and there thus remain the following:
\weq{2}{\begin{gathered}
Q = 1, \quad (0,1)(2,3), \quad (0,2)(1,3), \quad (0,3)(1,2) \\
(0,1,2), \quad (0,1,3), \quad (0,2,3), \quad (1,2,3) \\
(0,2,1), \quad (0,3,1), \quad (0,3,2), \quad (1,3,2).
\end{gathered}}

We now wish to seek out all the existing divisors of the group $P$. In doing so we leave aside the intransitive groups, which all consist either of three or of six permutations of only three elements, or of the permutations of two pairs of two elements. We therefore seek out only the \gesperrt{transitive divisors} of $P$. Moreover, in the case of a system of conjugate divisors we restrict ourselves to one representative, from which the others can indeed be derived by permutation of the figures (§. 154, 6.).

An essential simplification is brought about by the following two theorems:

1. \gesperrt{If a transitive permutation group of four figures contains a three-term cycle,}
%CONT>>
